% ============================================================
% meta.tex  --  SATU-SATUNYA tempat mengisi data diri & skripsi.
% Dipakai oleh halaman judul, pernyataan, pengesahan, intisari, abstract.
% ============================================================
\usepackage{graphicx}

\newcommand{\judul}{ANALISIS PENURUNAN TEKANAN DAN TRANSFER KALOR HEAT PIPE DI REAKTOR ANTARIKSA MENGGUNAKAN PYTHON}
\newcommand{\judulinggris}{Thesis Title in English}

\newcommand{\nama}{Muhammad Naufal Tamimi}
\newcommand{\nim}{20/456315/TK/50445}
\newcommand{\tahundaftar}{2020}
\newcommand{\tahun}{2026}

% Teknik Nuklir / Teknik Fisika
\newcommand{\prodi}{Teknik Nuklir}
% Nuclear Engineering / Engineering Physics
\newcommand{\prodien}{Nuclear Engineering}

\newcommand{\kotatanggal}{Yogyakarta, tanggal-bulan-tahun}
\newcommand{\tglujian}{tanggal bulan tahun ujian}
\newcommand{\tglpengesahan}{DD BULAN TAHUN}
\newcommand{\tglujianinggris}{Month Date, year}

\newcommand{\pembimbingutama}{Nama Lengkap Pembimbing Utama}
\newcommand{\pembimbingpendamping}{Nama Lengkap Pembimbing Pendamping}
\newcommand{\ketuasidang}{Nama Lengkap Ketua Sidang}
\newcommand{\pengujiutama}{Nama Lengkap Penguji Utama}
\newcommand{\anggotapenguji}{Nama Lengkap Anggota Penguji}

% Data Ketua Departemen
\newcommand{\ketuadepartemen}{Dr. Ir. Alexander Agung, S.T., M.Sc., IPU}
\newcommand{\nipketua}{19720916 199803 1002}

% Logo UGM berdiameter 5,5 cm. Simpan sebagai fig/logo-ugm.png
\newcommand{\logougm}{%
  \IfFileExists{fig/logo-ugm.png}%
    {\includegraphics[width=5.5cm]{fig/logo-ugm.png}}%
    {\fbox{\parbox[c][5.5cm][c]{5.5cm}{\centering\small Letakkan logo UGM\\(fig/logo-ugm.png)}}}%
}
